\documentclass{amsart}
% For dealing with references we use the comment environment
\usepackage{verbatim}
\newenvironment{reference}{\comment}{\endcomment}
%\newenvironment{reference}{}{}
\newenvironment{slogan}{\comment}{\endcomment}
\newenvironment{history}{\comment}{\endcomment}

% For commutative diagrams we use Xy-pic
\usepackage[all]{xy}

% We use 2cell for 2-commutative diagrams.
\xyoption{2cell}
\UseAllTwocells

% We use multicol for the list of chapters between chapters
\usepackage{multicol}

% This is generally recommended for better output
\usepackage{lmodern}
\usepackage[T1]{fontenc}

% For cross-file-references

% Package for hypertext links:
\usepackage{hyperref}

% For any local file, say "hello.tex" you want to link to please

% Theorem environments.
%
\theoremstyle{plain}
\newtheorem{theorem}[subsection]{Theorem}
\newtheorem{proposition}[subsection]{Proposition}
\newtheorem{lemma}[subsection]{Lemma}

\theoremstyle{definition}
\newtheorem{definition}[subsection]{Definition}
\newtheorem{example}[subsection]{Example}
\newtheorem{exercise}[subsection]{Exercise}
\newtheorem{situation}[subsection]{Situation}

\theoremstyle{remark}
\newtheorem{remark}[subsection]{Remark}
\newtheorem{remarks}[subsection]{Remarks}

\numberwithin{equation}{subsection}

% Macros
%
\def\lim{\mathop{\mathrm{lim}}\nolimits}
\def\colim{\mathop{\mathrm{colim}}\nolimits}
\def\Spec{\mathop{\mathrm{Spec}}}
\def\Hom{\mathop{\mathrm{Hom}}\nolimits}
\def\Ext{\mathop{\mathrm{Ext}}\nolimits}
\def\SheafHom{\mathop{\mathcal{H}\!\mathit{om}}\nolimits}
\def\SheafExt{\mathop{\mathcal{E}\!\mathit{xt}}\nolimits}
\def\Sch{\mathit{Sch}}
\def\Mor{\mathop{\mathrm{Mor}}\nolimits}
\def\Ob{\mathop{\mathrm{Ob}}\nolimits}
\def\Sh{\mathop{\mathit{Sh}}\nolimits}
\def\NL{\mathop{N\!L}\nolimits}
\def\CH{\mathop{\mathrm{CH}}\nolimits}
\def\proetale{{pro\text{-}\acute{e}tale}}
\def\etale{{\acute{e}tale}}
\def\QCoh{\mathit{QCoh}}
\def\Ker{\mathop{\mathrm{Ker}}}
\def\Im{\mathop{\mathrm{Im}}}
\def\Coker{\mathop{\mathrm{Coker}}}
\def\Coim{\mathop{\mathrm{Coim}}}

% Boxtimes
%
\DeclareMathSymbol{\boxtimes}{\mathbin}{AMSa}{"02}

%
% Macros for moduli stacks/spaces
%
\def\QCohstack{\mathcal{QC}\!\mathit{oh}}
\def\Cohstack{\mathcal{C}\!\mathit{oh}}
\def\Spacesstack{\mathcal{S}\!\mathit{paces}}
\def\Quotfunctor{\mathrm{Quot}}
\def\Hilbfunctor{\mathrm{Hilb}}
\def\Curvesstack{\mathcal{C}\!\mathit{urves}}
\def\Polarizedstack{\mathcal{P}\!\mathit{olarized}}
\def\Complexesstack{\mathcal{C}\!\mathit{omplexes}}
% \Pic is the operator that assigns to X its picard group, usage \Pic(X)
% \Picardstack_{X/B} denotes the Picard stack of X over B
% \Picardfunctor_{X/B} denotes the Picard functor of X over B
\def\Pic{\mathop{\mathrm{Pic}}\nolimits}
\def\Picardstack{\mathcal{P}\!\mathit{ic}}
\def\Picardfunctor{\mathrm{Pic}}
\def\Deformationcategory{\mathcal{D}\!\mathit{ef}}

\setlength{\emergencystretch}{3em}
\begin{document}
\title[Stacks Project, Tag 07NJ]{343 --- Continuous lifting through closed topologically nilpotent ideals}
\maketitle
\noindent Copyright (C) 2005--2025 Johan de Jong. Stacks Project authors. Source: \href{https://stacks.math.columbia.edu/tag/07NJ}{Tag 07NJ}.

\noindent Editorial difficulty note (not author text). Difficulty H2: The author constructs a compatible two-parameter grid of lifts through square-zero fibre-product surjections. A diagonal inverse limit produces the lift, and a quantitative power estimate proves continuity; ordinary formal smoothness alone does not directly supply these compatible topological lifts..

\noindent Editorial provenance note (not author text). History: excerpt from revision \texttt{a0444\allowbreak 6e57e\allowbreak c1fbc\allowbreak 252a8\allowbreak 71afc\allowbreak ec775\allowbreak 2fb28\allowbreak 07b14}, more-algebra.tex, lines 9256--9320. Reference presentation was adapted; original-author context and core support blocks were retained or added. The mathematical wording is unchanged. Licensed under GNU FDL 1.2 or later; no invariant sections or cover texts. See ../licenses/COPYING. Context blocks reproduce author definitions and supporting results; they are not additional counted problems.

\begin{lemma}
\label{lemma-lift-continuous}
Let $R \to S$ be a ring map. Let $\mathfrak n$ be an ideal of $S$.
Assume that $R \to S$ is formally smooth in the $\mathfrak n$-adic
topology. Consider a solid commutative diagram
$$
\xymatrix{
S \ar[r]_\psi \ar@{-->}[rd] & A/J \\
R \ar[r] \ar[u] & A \ar[u]
}
$$
of homomorphisms of topological rings where $A$ is adic
and $A/J$ is the quotient (as topological ring) of $A$ by a closed ideal
$J \subset A$ such that $J^t$ is contained in an ideal of definition
of $A$ for some $t \geq 1$. Then there exists a dotted arrow in the category of
topological rings which makes the diagram commute.
\end{lemma}

\begin{proof}
Let $I \subset A$ be an ideal of definition so that $I \supset J^t$
for some $t$. Then $A = \lim A/I^n$ and $A/J = \lim A/J + I^n$
because $J$ is assumed closed. Consider the following diagram of
discrete $R$ algebras $A_{n, m} = A/J^n + I^m$:
$$
\xymatrix{
A/J^3 + I^3 \ar[r] \ar[d] &
A/J^2 + I^3 \ar[r] \ar[d] &
A/J + I^3 \ar[d] \\
A/J^3 + I^2 \ar[r] \ar[d] &
A/J^2 + I^2 \ar[r] \ar[d] &
A/J + I^2 \ar[d] \\
A/J^3 + I \ar[r] &
A/J^2 + I \ar[r] &
A/J + I
}
$$
Note that each of the commutative squares defines a surjection
$$
A_{n + 1, m + 1} \longrightarrow A_{n + 1, m} \times_{A_{n, m}} A_{n, m + 1}
$$
of $R$-algebras whose kernel has square zero.
We will inductively construct $R$-algebra maps
$\varphi_{n, m} : S \to A_{n, m}$.
Namely, we have the maps $\varphi_{1, m} = \psi \bmod J + I^m$.
Note that each of these maps is continuous as $\psi$ is.
We can inductively choose the maps $\varphi_{n, 1}$ by starting
with our choice of $\varphi_{1, 1}$ and lifting up, using the
formal smoothness of $S$ over $R$, along the bottom row of the
diagram above. We construct the remaining maps $\varphi_{n, m}$
by induction on $n + m$. Namely, we choose $\varphi_{n + 1, m + 1}$
by lifting the pair $(\varphi_{n + 1, m}, \varphi_{n, m + 1})$
along the displayed surjection above (again using the formal smoothness
of $S$ over $R$). In this way all of the maps $\varphi_{n, m}$ are
compatible with the transition maps of the system.
As $J^t \subset I$ we see that for example
$\varphi_n = \varphi_{nt, n} \bmod I^n$ induces a map $S \to A/I^n$.
Taking the limit $\varphi = \lim \varphi_n$ we obtain a map
$S \to A = \lim A/I^n$. The composition into $A/J$ agrees
with $\psi$ as we have seen that $A/J = \lim A/J + I^n$.
Finally we show that $\varphi$ is continuous. Namely, we know that
$\psi(\mathfrak n^r) \subset J + I/J$ for some $r \geq 1$ by our
assumption that $\psi$ is a morphism of topological rings, see
Lemma \href{https://stacks.math.columbia.edu/tag/07E9}{lemma continuous (Tag 07E9)}. Hence $\varphi(\mathfrak n^r) \subset J + I$
hence $\varphi(\mathfrak n^{rt}) \subset I$ as desired.
\end{proof}
\end{document}
